%! Author = sriadityavedantam
%! Date = 3/30/26

\section{Chapter 5}

\begin{problem}{
    Which of the following subsets of $R[x]$ are subrings of $R[x]$?
    Justify your answer.

    (a) All polynomials with constant term $0_R$.

    (b) All polynomials of degree $2$.

    (c) All polynomials of degree $\leq k$, where $k$ is a fixed positive integer.

    (d) All polynomials in which the odd powers of $x$ have zero coefficients.

    (e) All polynomials in which the even powers of $x$ have zero coefficients.
}
    \textbf{(a).}
    Let
    \[
        S[x]\coloneqq \{f(x)\in R[x] : \text{the constant term of } f(x)\text{ is }0_R\}.
    \]
    We show that $S[x]$ is a subring of $R[x]$.
    If $y(x),z(x)\in S[x]$, then both have constant term $0_R$.
    Therefore the constant term of $y(x)+z(x)$ is
    \[
        0_R+0_R=0_R,
    \]
    so $y(x)+z(x)\in S[x]$.
    Also, the constant term of $-y(x)$ is
    \[
        -0_R=0_R,
    \]
    so $-y(x)\in S[x]$.
    For multiplication, if $y(x),z(x)\in S[x]$, then the constant term of $y(x)z(x)$ is the product of the constant terms:
    \[
        0_R\cdot 0_R=0_R.
    \]
    So $y(x)z(x)\in S[x]$.
    Finally, the zero polynomial is in $S[x]$.
    Thus $S[x]$ is a subring of $R[x]$.

    \textbf{(b).}
    This is not a subring of $R[x]$.
    A subring must contain the zero polynomial.
    But the zero polynomial does not have degree $2$.
    Also, the sum of two degree $2$ polynomials need not have degree $2$.
    For example, in any ring,
    \[
        x^2+1
        \quad\text{and}\quad
        -(x^2+1)
    \]
    both have degree $2$, but their sum is the zero polynomial.
    Thus this set is not a subring.

    \textbf{(c).}
    This is not a subring of $R[x]$ in general.
    It is closed under addition and contains the zero polynomial, but it is not closed under multiplication.
    If $f(x)$ and $g(x)$ both have degree $k$, then in general
    \[
        \deg(f(x)g(x))=2k,
    \]
    which is greater than $k$.
    Therefore this set is not a subring.

    \textbf{(d).}
    Let $S[x]\subseteq R[x]$ be the set of all polynomials in which the odd powers of $x$ have zero coefficients.
    Then every element of $S[x]$ can be written in the form
    \[
        \sum_{i=0}^n a_i x^{2i}.
    \]
    If
    \[
        y(x)=\sum_{i=0}^n a_i x^{2i}
        \qquad\text{and}\qquad
        z(x)=\sum_{j=0}^m b_j x^{2j},
    \]
    then
    \[
        y(x)+z(x)
    \]
    still has only even powers of $x$.
    So $S[x]$ is closed under addition.
    Also,
    \[
        -y(x)=\sum_{i=0}^n (-a_i)x^{2i},
    \]
    so $S[x]$ is closed under additive inverses.
    For multiplication,
    \begin{align*}
        y(x)z(x)
        &=
        \left(\sum_{i=0}^n a_i x^{2i}\right)\left(\sum_{j=0}^m b_j x^{2j}\right).
    \end{align*}
    Every term in this product has the form
    \[
        a_i b_j x^{2i+2j}=a_i b_j x^{2(i+j)},
    \]
    which again has even exponent.
    So $y(x)z(x)\in S[x]$.
    Also, the zero polynomial lies in $S[x]$.
    Thus $S[x]$ is a subring of $R[x]$.

    \textbf{(e).}
    This is not a subring.
    It is not closed under multiplication.
    For example,
    \[
        x\in S[x],
    \]
    but
    \[
        x\cdot x=x^2,
    \]
    and $x^2$ has a nonzero even-power coefficient.
    So $x^2\notin S[x]$.
    Also, the zero polynomial issue depends on convention, but the failure of closure under multiplication already shows this is not a subring.
\end{problem}


\begin{problem}{
    Use the Euclidean Algorithm to find the gcd of
    \[
        x^4+3x^3+2x+4
        \quad\text{and}\quad
        x^2-1
    \]
    in $\mathbb{Z}_5[x]$.
}
    Using the division algorithm in $\mathbb{Z}_5[x]$, divide
    \[
        x^4+3x^3+2x+4
    \]
    by
    \[
        x^2-1.
    \]
    First,
    \begin{align*}
        x^4+3x^3+2x+4
        &=
        x^2(x^2-1)+(3x^3+x^2+2x+4).
    \end{align*}
    Now divide the remainder by $x^2-1$ again:
    \begin{align*}
        3x^3+x^2+2x+4
        &=
        3x(x^2-1)+(x^2+5x+4)\\
        &=
        3x(x^2-1)+(x^2+4).
    \end{align*}
    Then
    \begin{align*}
        x^2+4
        &=
        1\cdot (x^2-1)+5\\
        &=
        1\cdot (x^2-1)+0.
    \end{align*}
    So altogether,
    \[
        x^4+3x^3+2x+4=(x^2+3x+1)(x^2-1)+0.
    \]
    Hence
    \[
        x^2-1\mid x^4+3x^3+2x+4.
    \]
    Therefore the gcd is
    \[
        \gcd(x^4+3x^3+2x+4,\ x^2-1)=x^2-1.
    \]
    Since $x^2-1$ is already monic, this is the gcd.
\end{problem}

\begin{problem}{
    Express the gcd from the previous problem as a linear combination of the two polynomials.
}
    Since
    \[
        \gcd(x^4+3x^3+2x+4,\ x^2-1)=x^2-1,
    \]
    we can write the gcd immediately as
    \[
        x^2-1
        =
        0\cdot (x^4+3x^3+2x+4)+1\cdot (x^2-1).
    \]
    Thus the gcd is a linear combination of the two polynomials.
\end{problem}

\begin{problem}{
    Let $f(x),g(x),h(x)\in F[x]$, with $f(x)$ and $g(x)$ relatively prime.
    If $f(x)\mid h(x)$ and $g(x)\mid h(x)$, prove that $f(x)g(x)\mid h(x)$.
}
    Since $f(x)$ and $g(x)$ are relatively prime, we have
    \[
        \gcd(f(x),g(x))=1.
    \]
    Therefore there exist $c(x),d(x)\in F[x]$ such that
    \[
        1=f(x)c(x)+g(x)d(x).
    \]
    Also, since $f(x)\mid h(x)$ and $g(x)\mid h(x)$, there exist $a(x),b(x)\in F[x]$ such that
    \[
        h(x)=f(x)a(x)
        \qquad\text{and}\qquad
        h(x)=g(x)b(x).
    \]
    Now multiply the Bézout identity by $h(x)$:
    \begin{align*}
        h(x)
        &=
        h(x)f(x)c(x)+h(x)g(x)d(x).
    \end{align*}
    Using the divisibility relations,
    \begin{align*}
        h(x)
        &=
        g(x)b(x)f(x)c(x)+f(x)a(x)g(x)d(x)\\
        &=
        f(x)g(x)\bigl(b(x)c(x)+a(x)d(x)\bigr).
    \end{align*}
    Therefore
    \[
        f(x)g(x)\mid h(x).
    \]
\end{problem}

\begin{problem}{
    (a) By counting products of the form $(x+a)(x+b)$, show that there are exactly
    \[
        \frac{p^2+p}{2}
    \]
    monic polynomials of degree $2$ that are not irreducible in $\mathbb{Z}_p[x]$.

    (b) Show that there are exactly
    \[
        \frac{p^2-p}{2}
    \]
    monic irreducible polynomials of degree $2$ in $\mathbb{Z}_p[x]$.
}
    \textbf{(a).}
    A monic quadratic in $\mathbb{Z}_p[x]$ has the form
    \[
        x^2+cx+d,
    \]
    where $c,d\in\mathbb{Z}_p$.
    So there are exactly
    \[
        p^2
    \]
    monic quadratics total.
    A monic quadratic is reducible if and only if it can be written as
    \[
        (x+a)(x+b)
    \]
    for some $a,b\in\mathbb{Z}_p$.
    Now count the number of distinct products of this form.
    Since
    \[
        (x+a)(x+b)=(x+b)(x+a),
    \]
    the ordered pairs $(a,b)$ and $(b,a)$ give the same polynomial.
    So we count unordered pairs $\{a,b\}$ with $a,b\in\mathbb{Z}_p$.
    The number of unordered pairs with repetition allowed from a set of size $p$ is
    \[
        \binom{p+1}{2}=\frac{p^2+p}{2}.
    \]
    Hence there are exactly
    \[
        \frac{p^2+p}{2}
    \]
    monic reducible polynomials of degree $2$ in $\mathbb{Z}_p[x]$.

    \textbf{(b).}
    Since every monic quadratic is either reducible or irreducible, the number of monic irreducible quadratics is
    \[
        p^2-\frac{p^2+p}{2}
        =
        \frac{2p^2-(p^2+p)}{2}
        =
        \frac{p^2-p}{2}.
    \]
    Thus there are exactly
    \[
        \frac{p^2-p}{2}
    \]
    monic irreducible polynomials of degree $2$ in $\mathbb{Z}_p[x]$.
\end{problem}

\begin{problem}{
    (a) Show that $x^2+2$ is irreducible in $\mathbb{Z}_5[x]$.

    (b) Factor $x^4-4$ as a product of irreducibles in $\mathbb{Z}_5[x]$.
}
    \textbf{(a).}
    A quadratic polynomial over a field is reducible if and only if it has a root.
    So we check whether
    \[
        x^2+2
    \]
    has a root in $\mathbb{Z}_5$.
    The squares in $\mathbb{Z}_5$ are
    \[
        0^2=0,\quad 1^2=1,\quad 2^2=4,\quad 3^2=4,\quad 4^2=1.
    \]
    So the only square classes are $0,1,4$.
    But
    \[
        x^2+2=0
        \iff
        x^2= -2 \equiv 3 \mod 5.
    \]
    Since $3$ is not a square in $\mathbb{Z}_5$, the polynomial has no root.
    Therefore
    \[
        x^2+2
    \]
    is irreducible in $\mathbb{Z}_5[x]$.

    \textbf{(b).}
    In $\mathbb{Z}_5[x]$, we have
    \[
        -4\equiv 1,
    \]
    so
    \[
        x^4-4=x^4+1.
    \]
    Now factor:
    \begin{align*}
        x^4+1
        &=
        (x^2+2)(x^2-2),
    \end{align*}
    since in $\mathbb{Z}_5$,
    \[
        (x^2+2)(x^2-2)=x^4-4=x^4+1.
    \]
    Now
    \[
        x^2+2
    \]
    is irreducible by part (a).
    Also,
    \[
        x^2-2=x^2+3.
    \]
    To see whether this is reducible, check whether $2$ is a square mod $5$.
    It is not, since the square classes are $0,1,4$.
    So $x^2-2$ has no root in $\mathbb{Z}_5$, and hence is irreducible.
    Therefore
    \[
        x^4-4=(x^2+2)(x^2-2)
    \]
    is the factorization into irreducibles in $\mathbb{Z}_5[x]$.
\end{problem}


\begin{problem}{
    Let $\phi:\mathbb{C}\mapsto\mathbb{C}$ be an isomorphism of rings such that $\phi(a)=a$ for each $a\in\mathbb{Q}$.
    Suppose $r\in\mathbb{C}$ is a root of $f(x)\in\mathbb{Q}[x]$.
    Prove that $\phi(r)$ is also a root of $f(x)$.
}
    Since $r$ is a root of $f(x)$, we have
    \[
        f(r)=0.
    \]
    Write
    \[
        f(x)=a_{n}x^n+a_{n-1}x^{n-1}+\cdots+a_{1}x+a_0,
    \]
    where each $a_i\in\mathbb{Q}$.
    Now evaluate $f$ at $\phi(r)$:
    \begin{align*}
        f(\phi(r))
        &=
        a_n(\phi(r))^n+a_{n-1}(\phi(r))^{n-1}+\cdots+a_1\phi(r)+a_0.
    \end{align*}
    Since $\phi$ is a ring isomorphism and fixes every rational number, we have
    \[
        \phi(a_i)=a_i
    \]
    for all $i$, and also
    \[
        \phi(r^k)=\phi(r)^k.
    \]
    Thus
    \begin{align*}
        f(\phi(r))
        &=
        \phi(a_n)\phi(r^n)+\phi(a_{n-1})\phi(r^{n-1})+\cdots+\phi(a_1)\phi(r)+\phi(a_0)\\
        &=
        \phi(a_{n}r^n+a_{n-1}r^{n-1}+\cdots+a_{1}r+a_0)\\
        &=
        \phi(f(r))\\
        &=
        \phi(0)\\
        &=
        0.
    \end{align*}
    Therefore $\phi(r)$ is also a root of $f(x)$.
\end{problem}

\begin{problem}{
    We say that $a\in F$ is a multiple root of $f(x)\in F[x]$ if $(x-a)^k$ is a factor of $f(x)$ for some $k\geq 2$.

    (a) Prove that $a\in\mathbb{R}$ is a multiple root of $f(x)\in\mathbb{R}[x]$ if and only if $a$ is a root of both $f(x)$ and $f'(x)$.

    (b) If $f(x)\in\mathbb{R}[x]$ and if $f(x)$ is relatively prime to $f'(x)$, prove that $f(x)$ has no multiple root in $\mathbb{R}[x]$.
}
    \textbf{(a).}
    Suppose first that $a$ is a multiple root of $f(x)$.
    Then
    \[
        f(x)=(x-a)^k g(x)
    \]
    for some $k\geq 2$ and some $g(x)\in\mathbb{R}[x]$.
    In particular, $(x-a)^2$ divides $f(x)$, so we may write
    \[
        f(x)=(x-a)^{2}g(x).
    \]
    Then clearly
    \[
        f(a)=0.
    \]
    Differentiate:
    \begin{align*}
        f'(x)
        &=
        2(x-a)g(x)+(x-a)^{2}g'(x)\\
        &=
        (x-a)\bigl(2g(x)+(x-a)g'(x)\bigr).
    \end{align*}
    Thus $(x-a)\mid f'(x)$, so
    \[
        f'(a)=0.
    \]
    Conversely, suppose $a$ is a root of both $f(x)$ and $f'(x)$.
    Since $f(a)=0$, we can write
    \[
        f(x)=(x-a)h(x)
    \]
    for some $h(x)\in\mathbb{R}[x]$.
    Differentiate:
    \[
        f'(x)=h(x)+(x-a)h'(x).
    \]
    Now evaluate at $x=a$:
    \[
        f'(a)=h(a).
    \]
    But $f'(a)=0$, so
    \[
        h(a)=0.
    \]
    Hence $(x-a)\mid h(x)$, so
    \[
        h(x)=(x-a)j(x)
    \]
    for some $j(x)\in\mathbb{R}[x]$.
    Therefore
    \[
        f(x)=(x-a)^{2}j(x),
    \]
    so $a$ is a multiple root of $f(x)$.

    \textbf{(b).}
    If $f(x)$ had a multiple root $a\in\mathbb{R}$, then by part (a), $a$ would be a root of both $f(x)$ and $f'(x)$.
    That would mean $(x-a)$ is a common factor of $f(x)$ and $f'(x)$.
    But this contradicts the assumption that $f(x)$ and $f'(x)$ are relatively prime.
    Therefore $f(x)$ has no multiple root in $\mathbb{R}[x]$.
\end{problem}

\begin{problem}{
    Factor $x^{12}-1$ over each of $\mathbb{R}[x]$ and $\mathbb{Q}[x]$.
}
    We first factor over $\mathbb{Q}[x]$:
    \begin{align*}
        x^{12}-1
        &=
        (x^6-1)(x^6+1)\\
        &=
        (x^3-1)(x^3+1)(x^6+1)\\
        &=
        (x-1)(x^2+x+1)(x+1)(x^2-x+1)(x^6+1).
    \end{align*}
    Now factor
    \[
        x^6+1=(x^2+1)(x^4-x^2+1).
    \]
    So over $\mathbb{Q}[x]$ we get
    \[
        x^{12}-1
        =
        (x-1)(x+1)(x^2+x+1)(x^2-x+1)(x^2+1)(x^4-x^2+1).
    \]
    Now factor further over $\mathbb{R}[x]$.
    The quartic
    \[
        x^4-x^2+1
    \]
    factors as
    \[
        x^4-x^2+1=\left(x^2+\sqrt{3}x+1\right)\left(x^2-\sqrt{3}x+1\right).
    \]
    Thus over $\mathbb{R}[x]$,
    \[
        x^{12}-1
        =
        (x-1)(x+1)(x^2+x+1)(x^2-x+1)(x^2+1)\left(x^2+\sqrt{3}x+1\right)\left(x^2-\sqrt{3}x+1\right).
    \]
    Therefore,
    \begin{align*}
        \mathbb{Q}[x]\ni x^{12}-1
        &=
        (x-1)(x+1)(x^2+x+1)(x^2-x+1)(x^2+1)(x^4-x^2+1),\\
        \mathbb{R}[x]\ni x^{12}-1
        &=
        (x-1)(x+1)(x^2+x+1)(x^2-x+1)(x^2+1)(x^2+\sqrt{3}x+1)(x^2-\sqrt{3}x+1).
    \end{align*}
\end{problem}

\begin{problem}{
    Prove that
    \[
        f(x)=x^4+x^3+x^2+x+1
    \]
    is irreducible in $\mathbb{Q}[x]$.
}
    We use Eisenstein's criterion after shifting by $x\mapsto x+1$.
    First note that
    \[
        (x-1)f(x)=x^5-1,
    \]
    so
    \[
        f(x)=\frac{x^5-1}{x-1}.
    \]
    Then
    \begin{align*}
        f(x+1)
        &=
        \frac{(x+1)^5-1}{(x+1)-1}\\
        &=
        \frac{(x+1)^5-1}{x}.
    \end{align*}
    Expand:
    \begin{align*}
        (x+1)^5-1
        &=
        x^5+5x^4+10x^3+10x^2+5x+1-1\\
        &=
        x^5+5x^4+10x^3+10x^2+5x.
    \end{align*}
    So
    \[
        f(x+1)=x^4+5x^3+10x^2+10x+5.
    \]
    Now apply Eisenstein's criterion with $p=5$.
    We have
    \[
        5\mid 5,\ 10,\ 10,\ 5,
    \]
    so $5$ divides every coefficient except the leading coefficient.
    Also,
    \[
        5\nmid 1,
    \]
    and
    \[
        25\nmid 5.
    \]
    Therefore $f(x+1)$ is irreducible in $\mathbb{Q}[x]$.
    A polynomial is irreducible if and only if its translate by $x\mapsto x+1$ is irreducible.
    Hence $f(x)$ is irreducible in $\mathbb{Q}[x]$.
\end{problem}

\begin{problem}{
    Prove that for $p$ prime,
    \[
        f(x)=x^{p-1}+x^{p-2}+\dots+x^2+x+1
    \]
    is irreducible in $\mathbb{Q}[x]$.
}
    We use Eisenstein's criterion after the substitution $x\mapsto x+1$.
    First note that
    \[
        (x-1)f(x)=x^p-1,
    \]
    so
    \[
        f(x)=\frac{x^p-1}{x-1}.
    \]
    Then
    \begin{align*}
        f(x+1)
        &=
        \frac{(x+1)^p-1}{x}.
    \end{align*}
    Using the binomial theorem,
    \begin{align*}
        (x+1)^p-1
        &=
        \sum_{n=0}^p \binom{p}{n}x^n-1\\
        &=
        \sum_{n=1}^p \binom{p}{n}x^n.
    \end{align*}
    So
    \begin{align*}
        f(x+1)
        &=
        \frac{\binom{p}{1}x+\binom{p}{2}x^2+\dots+\binom{p}{p-1}x^{p-1}+\binom{p}{p}x^p}{x}\\
        &=
        \binom{p}{1}+\binom{p}{2}x+\dots+\binom{p}{p-1}x^{p-2}+\binom{p}{p}x^{p-1}.
    \end{align*}
    Now for $1\leq k\leq p-1$, we know that
    \[
        p\mid \binom{p}{k}.
    \]
    Also,
    \[
        \binom{p}{p}=1,
    \]
    so $p$ does not divide the leading coefficient.
    Finally,
    \[
        \binom{p}{1}=p,
    \]
    and
    \[
        p^2\nmid p.
    \]
    Thus Eisenstein's criterion applies with the prime $p$.
    Therefore $f(x+1)$ is irreducible in $\mathbb{Q}[x]$.
    Hence $f(x)$ is irreducible in $\mathbb{Q}[x]$.
\end{problem}